\documentclass[11pt]{article}

\usepackage[T1]{fontenc}
\usepackage{lmodern}
\usepackage[margin=1.02in]{geometry}
\usepackage{amsmath,amssymb,amsthm,mathtools}
\usepackage{microtype}
\usepackage{graphicx}
\usepackage[hidelinks]{hyperref}
\usepackage{enumitem}

\setlist{nosep}
\newtheorem{theorem}{Theorem}
\newtheorem{lemma}{Lemma}
\theoremstyle{remark}
\newtheorem{remark}{Remark}

\newcommand{\C}{\mathbb C}
\newcommand{\Sph}{\mathbb S}
\newcommand{\cL}{\mathcal L}
\newcommand{\cK}{\mathcal K}
\newcommand{\Cal}{\operatorname{Cal}}
\newcommand{\Exp}{\operatorname{Exp}}
\newcommand{\Inv}{\operatorname{Inv}}
\newcommand{\eps}{\epsilon}

\title{Noncommuting Exponential Spectra in an Operator Algebra}
\author{Independent research note}
\date{August 17, 2026}

\begin{document}
\maketitle

\begin{abstract}
Klaja and Ransford asked whether there are a Banach space $E$ and bounded
operators $S,T$ on $E$ for which the nonzero exponential spectra of $ST$
and $TS$ differ.  We give a positive answer, with $E$ separable.  The
construction combines their explicit matrix functions on $\Sph^4$ with
Motakis's Banach space whose Calkin algebra is $C(\Sph^4)$.  The crucial
extra feature is that every Lipschitz scalar function has an exact
multiplicative bounded lift.  Matrix amplification lifts the
Klaja--Ransford pair: the Calkin quotient detects that one product is not in
the identity component, while the reversed product is a single explicit
operator exponential.
\end{abstract}

\section*{Packet status}

\begin{description}[leftmargin=3.1cm,style=nextline]
\item[Run] \texttt{fa\_banach\_001}, lane 7.
\item[Source] H.~Klaja and T.~Ransford,
\emph{Non-commutativity of the exponential spectrum}, arXiv:1510.08109,
Question~4.2, PDF page~7 \cite{KR}.
\item[Classification] Candidate full result, likely valid subject to human
mathematical review.
\item[Novelty] Moderate confidence.  The proof is a new synthesis of two
published constructions; no prior connection was located in a bounded
search through August 17, 2026.
\end{description}

\begin{center}
\includegraphics[width=0.94\textwidth]{figures/open_problem_crop.png}
\end{center}

For a unital complex Banach algebra $A$, let $\Exp(A)$ denote the connected
component of $1$ in $\Inv(A)$, equivalently the finite products of
exponentials, and put
\[
 \eps(x)=\{\lambda\in\C:\lambda1-x\notin\Exp(A)\}.
\]
The source asks whether $A=\cL(E)$ can contain $S,T$ with
$\eps(ST)\setminus\{0\}\ne\eps(TS)\setminus\{0\}$.

\section{Proof intuition}

Klaja and Ransford already construct smooth $2\times2$ matrix functions
$a,b$ on $\Sph^4$ for which $1-2ab$ is invertible but not null-homotopic,
whereas $1-2ba$ is null-homotopic.  A representation of this function
algebra by operators is not enough: a path may appear after embedding into
the much larger algebra $\cL(E)$.

The remedy is to combine a lift with a quotient detector.  Motakis
constructs a separable space $X$ with $\Cal(X)\cong C(\Sph^4)$ and bounded
diagonal lifts $f\mapsto D_f$ for every Lipschitz $f$; on Lipschitz
functions this lift is exactly multiplicative.  On $E=X\oplus X$, entrywise
amplification lifts the smooth matrix functions, while the Calkin quotient
maps them back to the originals.  Thus the hard component cannot disappear.
For the easy component we do not need to lift an arbitrary homotopy: its
scalar entry has the explicit Lipschitz logarithm
$i(\pi-4\arctan t)$, so its operator lift is one exponential.

\section{The Lipschitz Calkin lift}

We use the following consequence of Motakis's Propositions~4.1 and~4.3 and
Corollary~5.3 \cite{Motakis}.

\begin{lemma}[Motakis lift]\label{lem:lift}
Let $K$ be a compact metric space.  There are a separable complex Banach
space $X$, a Schauder basis $(d_\gamma)_{\gamma\in\Gamma}$, and a map
$\kappa:\Gamma\to K$ such that:
\begin{enumerate}[label=\textup{(\roman*)},leftmargin=1.05cm]
\item for every scalar Lipschitz function $f$ on $K$, the rule
\[
 D_f d_\gamma=f(\kappa(\gamma))d_\gamma
\]
extends to a bounded operator on $X$;
\item $f\mapsto D_f$ is a unital algebra homomorphism on
$\operatorname{Lip}(K)$;
\item if $q_X:\cL(X)\to\Cal(X)$ is the quotient map, then
$f\mapsto q_X(D_f)$ extends to a Banach algebra isomorphism
$\Psi:C(K)\to\Cal(X)$.
\end{enumerate}
\end{lemma}

Take $K=\Sph^4$ and this space $X$, and put $E=X\oplus X$ (with, for
example, the maximum norm).  Block matrices give canonical Banach algebra
identifications
\[
 \cL(E)=M_2(\cL(X)),\qquad
 \cK(E)=M_2(\cK(X)),\qquad
 \Cal(E)=M_2(\Cal(X)).                                      \tag{1}
\]
The compact identity follows because a finite block operator is compact if
and only if each entry is compact.

For a matrix-valued function $F=(f_{ij})$ whose entries are Lipschitz, define
\[
 J(F):=(D_{f_{ij}})_{i,j=1}^2\in\cL(E).
\]
Lemma~\ref{lem:lift}(ii) implies
\[
 J(FG)=J(F)J(G),\qquad J(I_2)=I_E.                           \tag{2}
\]
Under (1), matrix amplification of $\Psi^{-1}$ defines an isomorphism
\[
 \Theta:\Cal(E)\longrightarrow M_2(C(K))
\]
such that
\[
 \Theta(q_E(J(F)))=F.                                       \tag{3}
\]
If both $F$ and $F^{-1}$ have Lipschitz entries, then (2) also gives
$J(F)^{-1}=J(F^{-1})$.

\section{The operator counterexample}

Write a point of $\Sph^4$ as $(z_0,z_1,t)\in\C^2\times\mathbb R$, where
$|z_0|^2+|z_1|^2+t^2=1$.  Define the smooth matrix functions
\[
 a(z_0,z_1,t)=\frac{1}{1+it}
 \begin{pmatrix}z_0&0\\z_1&0\end{pmatrix},\qquad
 b(z_0,z_1,t)=\frac{1}{1+it}
 \begin{pmatrix}\overline z_0&\overline z_1\\0&0\end{pmatrix}.       \tag{4}
\]
Klaja and Ransford prove that
\[
 c:=I_2-2ab:\Sph^4\longrightarrow GL_2(\C)                 \tag{5}
\]
is not null-homotopic \cite[Theorems~1.1 and~1.2]{KR}.  Equivalently,
\[
 c\notin\Exp(C(\Sph^4,M_2(\C))).                           \tag{6}
\]
All entries of $a,b,c$ are Lipschitz.  Since $c$ is a smooth map into
$GL_2(\C)$ on a compact manifold, $c^{-1}$ also has Lipschitz entries.

\begin{theorem}\label{thm:main}
There exist a separable complex Banach space $E$ and operators
$S,T\in\cL(E)$ such that
\[
 \frac12\in\eps(ST)\quad\text{and}\quad
 \frac12\notin\eps(TS).
\]
Consequently Question~4.2 of Klaja--Ransford has a positive answer.
\end{theorem}

\begin{proof}
With $E=X\oplus X$ as above, set
\[
 S=J(a),\qquad T=J(b).
\]
Equations (2) and (5) give
\[
 I_E-2ST=J(c).                                               \tag{7}
\]
This operator is invertible, with inverse $J(c^{-1})$.  If it belonged to
$\Exp(\cL(E))$, then applying the unital quotient homomorphism $q_E$ and the
isomorphism $\Theta$ would, by (3), put $c$ in
$\Exp(C(\Sph^4,M_2(\C)))$, contradicting (6).  Hence
\[
 I_E-2ST\notin\Exp(\cL(E)).                                 \tag{8}
\]

For the reversed product, direct multiplication in (4) gives
\[
 I_2-2ba=\begin{pmatrix}\phi(t)&0\\0&1\end{pmatrix},\qquad
 \phi(t)=-\left(\frac{1-it}{1+it}\right)^2.                \tag{9}
\]
For $-1\le t\le1$, put
\[
 g(t)=i\bigl(\pi-4\arctan t\bigr).
\]
The identity
\[
 \frac{1-it}{1+it}=e^{-2i\arctan t}
\]
shows that $\phi=e^g$.  In particular $g$ and $\phi$ are Lipschitz.  If
$G=\operatorname{diag}(g,0)$, then (2), or simply evaluation on every basis
vector, yields
\[
 I_E-2TS
 =J(\operatorname{diag}(\phi,1))
 =\exp(J(G))\in\Exp(\cL(E)).                                \tag{10}
\]

Finally,
\[
 \tfrac12I_E-ST=\tfrac12(I_E-2ST),\qquad
 \tfrac12I_E-TS=\tfrac12(I_E-2TS).
\]
Since $\tfrac12I_E=\exp(-\log(2)I_E)$ belongs to
$\Exp(\cL(E))$ and $\Exp(\cL(E))$ is a subgroup, multiplication by this
scalar preserves membership in $\Exp(\cL(E))$.  Statements (8) and (10)
therefore give the two asserted exponential-spectrum memberships.
\end{proof}

\section{Verification, scope, and novelty}

The proof uses only exact algebraic identities after the two cited
constructions.  The finite direct-sum Calkin identity (1), multiplicativity
(2), quotient detection (3), regularity of $c^{-1}$, the orientation of
$ST$ versus $TS$, and the logarithm in (9)--(10) were audited separately.
No numerical computation is used.

The result settles Question~4.2 as stated and strengthens it by making $E$
separable.  It does not settle the adjacent Question~4.1, which asks for a
universal torsion consequence for arbitrary Banach algebras.  The space $E$
is obtained from a Bourgain--Delbaen construction and is not claimed to be a
classical sequence or function space.

A bounded novelty search covered the run indexes, the local arXiv corpus,
exact question text, author/title combinations, Calkin realizations, and a
citation list for \cite{KR}.  The closest later result found is
Daniel--Ghosh \cite{DG}, which proves commutativity in
$\cL(\ell^p\oplus\ell^q)$ and states that Calkin commutativity implies
operator-algebra commutativity.  It does not give the existence theorem
above and does not cite Motakis.  No source connecting \cite{KR} and
\cite{Motakis} was located.  Because the ingredients are published even
though their combination appears new, novelty confidence is moderate.

Human review should focus on (1)--(3) and the explicit logarithm in (10).

\begin{thebibliography}{9}
\bibitem{KR}
H.~Klaja and T.~Ransford,
\emph{Non-commutativity of the exponential spectrum},
J. Funct. Anal. 272 (2017), 4158--4164; arXiv:1510.08109.

\bibitem{Motakis}
P.~Motakis,
\emph{Separable Spaces of Continuous Functions as Calkin Algebras},
arXiv:2110.10868.

\bibitem{DG}
S.~Daniel and A.~Ghosh,
\emph{A note on commutativity of the exponential spectrum in the operator
algebra $\mathcal B(\ell^p\oplus\ell^q)$},
Acta Math. Hungar. 177 (2025), 280--296,
doi:10.1007/s10474-025-01571-x.
\end{thebibliography}

\end{document}
